%% TODO -- confirm number of detectors; currently six (6)


We present a new method for finding reactor neutrino angular directionality developed via simulation across six low-energy electron-antineutrino detector designs. We highlight particularly interesting results as part of our larger project with the
goal of optimising their ability to resolve the neutrino source direction. Herein we focus upon the ability of various conceptual designs not only to detect the Inverse Beta Decay reaction, but to determine the origin direction of the neutrinos.
While this has secondary utility for single reactor power observation, it can yield important information about the burning profile and constrain the clandestine production of weapons material. In a multi-reactor scenario, this methodology has the capability to determine the products of individual reactors. 



%The structure of this paper is as follows: 
% \hyperref[sec:intro]{Section I} explains the problem at hand and serves as an introduction. \hyperref[sec:dirn]{Section II} covers how to find directionality in the context of Inverse Beta Decay. \hyperref[sec:detectors]{Section III} covers detector geometries and hybrid setups. \hyperref[sec:method]{Section IV} is methodology, including simulations and addressing backgrounds. \hyperref[sec:results]{Section V} presents our results. \hyperref[sec:newdesign]{Section VI}  considers new designs. \hyperref[sec:deltaphi]{Section VII} discusses our novel algorithm for determining direction and resolution. In \hyperref[sec:future]{Section VIII} we consider future directions of work, and we overview and conclude this paper in \hyperref[sec:concl]{Section IX}.


% \S\ref{sec:intro} explains the problem at hand and serves as an introduction.
% \S\ref{sec:dirn} covers how to find directionality in the context of Inverse Beta Decay.
% \S\ref{sec:detectors} covers detector geometries and hybrid setups.
% \S\ref{sec:method} is methodology, including simulations and addressing backgrounds.
% \S\ref{sec:results} presents our results.
% \S\ref{sec:newdesign}  considers new designs.
% \S\ref{sec:deltaphi} discusses our novel algorithm for determining direction and resolution.
% In \S\ref{sec:future} we consider future directions of work, and we overview and conclude this paper in \S\ref{sec:concl}.

